\section{An independent certificate for the polynomial}\label{sec:certificate}

Starting from the coefficients of $f$, we determine its discriminant,
remove the apparent branch points, and certify two monodromy permutations.
Their group then gives both geometric monodromy and regularity.
Put $q(t)=t^2+t+1$.

\begin{proposition}[Computational certificate]\label{prop:certificate}
The following statements hold for $f$.
\begin{enumerate}[label=\textup{(\roman*)},leftmargin=*]
\item Its exact $X$-discriminant is
\begin{equation}\label{eq:discriminant}
 \disc_X(f)=c\,q(t)^{40}S(t)^2,
\end{equation}
where $c\in\Q^\times$, $S\in\Q[t]$ is monic and squarefree of degree
$342$, and $\gcd(S,q)=1$.
\item At each root $a$ of $S$, $f(X,a)$ has exactly one repeated root,
and that root has multiplicity exactly $2$.
\item Over $\Q(\zetaT)$ one has
$f(X,\zetaT)=f_5(X)f_{20}(X)^3$, with distinct irreducible factors of the
indicated degrees, and
$\Res_X(f_{20},f_t(X,\zetaT))\ne0$.
\item Continuation from $t=0$ along standard meridians around $\zetaT$
and $\zetaT^2$ gives permutations $\sigma_1,\sigma_2$ of type
$1^5 3^{20}$. With permutations composed from left to right,
$\sigma_\infty=(\sigma_1\sigma_2)^{-1}$ has type $13^5$, and they generate a
transitive primitive group of order $87360$.
\end{enumerate}
\end{proposition}

The following computations establish the proposition; the programs and
exact outputs are documented in \cref{sec:computation}.

\subsection{The exact discriminant}

The general discriminant degree bound is
$\deg_t\disc_X(f)\leq2(65-1)7=896$.
Exact evaluations at $897$ distinct rational values of $t$ and exact
interpolation therefore determine the discriminant. Additional evaluations
check the interpolation.
Exact division and gcd computations give \eqref{eq:discriminant},
squarefreeness, and coprimality. The resulting discriminant has degree
$764=80+2\cdot342$.

\subsection{Nodes of the plane model}

\begin{lemma}\label{lem:node}
Let $F(X,t)$ be monic in $X$ with coefficients holomorphic near $t=a$.
Suppose $\ord_{t=a}\disc_X(F)=2$ and $F(X,a)$ has exactly one repeated
root, of multiplicity $2$. Then the normalization of $F(X,t)=0$ is
unramified above $a$.
\end{lemma}
\begin{proof}
Separate the simple roots holomorphically. Weierstrass preparation at the
double root reduces the remaining two roots to a quadratic factor
$X^2+u(t)X+v(t)$, after translating $X$. Its discriminant has order exactly
$2$, because the discriminants and mutual resultants of the other factors
are units. Thus
$u(t)^2-4v(t)=(t-a)^2w(t)$ with $w(a)\ne0$. A local holomorphic square root
of $w$ gives the two roots
\[
 \frac{-u(t)\pm(t-a)\sqrt{w(t)}}{2}.
\]
They are single-valued holomorphic functions of $t$.
\end{proof}

There is an exact modular check of \cref{prop:certificate}\textup{(ii)}.
Let $B(t)$ be the first principal subresultant coefficient of $f$ and
$f_X$, considered as polynomials in $X$. At a root $a$ of $S$, their
resultant vanishes, and $B(a)\ne0$ is equivalent to
$\deg\gcd(f(X,a),f_X(X,a))=1$. In characteristic zero this degree is
$\sum_j(m_j-1)$, where $m_j$ are the root multiplicities. Degree one
therefore means exactly one double root, with all other roots simple.

It suffices to prove $\Res_t(S,B)\ne0$ modulo one suitable prime.
At $p=1000003$, exact discriminant interpolation modulo $p$ gives degree
$764$ and the shape $\bar c\,\bar q^{40}\bar S^2$, with $\bar S$ monic
and squarefree of degree $342$ and coprime to $\bar q$.
The full-degree reduction ensures that the characteristic-zero monic
factor $S$ is $p$-integral and reduces to this $\bar S$.
Its irreducible factors have degrees
\[
 1,\ 3,\ 5,\ 14,\ 78,\ 84,\ 157.
\]
For every factor $P$, the gcd of $\bar f$ and $\bar f_X$ over
$\mathbb F_p[t]/(P)$ has degree one. Subresultants commute with these
specializations because the leading coefficients $1$ and $65$ remain
units. Hence $\bar B$ is nonzero at every root of $\bar S$, proving the
required nonvanishing in characteristic zero. The independent verifier
\texttt{scripts/check\_nodes\_modp.py} runs this finite-field calculation
from the coefficients of $f$; it uses the exact discriminant identity
\eqref{eq:discriminant} to draw its characteristic-zero conclusion.

The certificate also contains an interval proof. For each root
of $S$, it first certifies a parameter rectangle
$T$ containing exactly one root by interval Newton on $S$. The $342$
rectangles are pairwise disjoint, hence account for all roots of $S$.
It then isolates all $64$ roots of $f_X(X,t)$ uniformly for $t\in T$ in
disjoint rectangles. The value of $f$ excludes zero on all but one of
these rectangles, and $f_{XX}$ excludes zero on the remaining rectangle.
The vanishing of the discriminant supplies existence of a repeated root;
these tests establish its uniqueness and multiplicity.

The enclosure conventions are recorded in \cref{sec:implementation}.
Together with \cref{lem:node}, either test proves that the normalized cover has
no finite branch point outside $\{\zetaT,\zetaT^2\}$.

At a root of $f_{20}$ over $t=\zeta_3$, the nonvanishing of $f_t$ in
\cref{prop:certificate}\textup{(iii)} makes the plane curve smooth and
expresses $t-\zeta_3$ as a function with a zero of order $3$ in a local
$X$-coordinate. The twenty such points account for discriminant order
$20(3-1)=40$. The five simple roots are unramified. Complex conjugation
gives the same description over $\zeta_3^2$.

\subsection{Certified continuation}

The base point is $t=0$. A \emph{standard meridian} goes from the base
point along a tail to a small neighborhood of one branch point, makes one
counter-clockwise circuit, and returns along the same tail. Here the tails
are straight, with entry point $z-\varepsilon z/|z|$ for
$z=\zeta_3,\zeta_3^2$; the circuits are inscribed polygons in disjoint
small neighborhoods of the branch points. Path parameters are recorded
in \cref{sec:implementation}.
The two meridians wind once about their respective branch point and zero
times about the other, and form a generating system for the thrice-punctured
sphere. Apparent branch points have already been removed by \cref{lem:node}.

At a parameter value $t_0$, isolate the roots of $f(X,t_0)$ in rectangles
with centers $c_i$ and enclosing Euclidean radii $r_i$.
Choose disjoint discs $|X-c_i|<R_i$ containing those rectangles.
On the $i$th boundary circle, the monic factorization gives
\begin{equation}\label{eq:lower}
 |f(X,t_0)|\geq
 (R_i-r_i)\prod_{\ell\ne i}(|c_i-c_\ell|-R_i-r_\ell).
\end{equation}
For each $1\leq j\leq7$, isolate the roots of $\partial_t^j f(X,t_0)$.
If its leading coefficient is $a_j$ and its roots have centers $b_{j\ell}$,
enclosing radii $\rho_{j\ell}$, and multiplicities $m_{j\ell}$, then
\[
 |\partial_t^j f(X,t_0)|\leq
 |a_j|\prod_\ell
 (|c_i-b_{j\ell}|+R_i+\rho_{j\ell})^{m_{j\ell}}.
\]
The Taylor expansion in $t$ terminates at degree $7$. Bounding its $j$th
term by $h^j/j!$ times this product gives a Rouch\'e test valid for every
$|t-t_0|\leq h$. For each candidate radius, bisection selects $h$ such
that the sum of the seven upper bounds is strictly below \eqref{eq:lower};
the common step is the minimum over all roots. Factoring the
$t$-derivatives avoids severe cancellation losses in bounds obtained from
absolute values of coefficients.

Each accepted endpoint is checked in ball arithmetic to lie within the
certified parameter displacement. Isolated endpoint roots are matched by
containment in the preceding discs. Thus labels persist along the whole
path. The resulting permutations are supplied in
\texttt{data/monodromy.json}.

There is a useful symmetry check. Let $c$ be complex conjugation on the
isolated roots of $f(X,0)$. The two meridians are conjugate up to homotopy
with reversed orientation, so
\[
 \sigma_2=c\,\sigma_1^{-1}c.
\]
Matching conjugate root enclosures gives an involution of type $1\,2^{32}$,
and the permutations obtained by tracking both loops satisfy this
identity exactly.

\subsection{Geometric group and regularity}

By \cref{lem:node}, the two meridians generate the fundamental group
needed for the normalized cover. Transitivity of their action proves
connectedness. The GAP primitive-group library has $13$ groups of degree
$65$; exactly one has order $87360$, namely $G=\Sz(8)\rtimes C_3$.
This identifies the geometric monodromy in
\cref{prop:certificate}\textup{(iv)}.

The arithmetic monodromy contains $G$ normally and is contained in
$N_{S_{65}}(G)$. This normalizer is $G$ itself. One can see this
without a second search: the simple derived subgroup $N$ is characteristic
in $G$, the centralizer of its nonregular doubly transitive action is
trivial, and $\Aut(N)=G$. Conjugation therefore embeds the normalizer in
$G$, while it already contains $G$. Arithmetic and geometric monodromy
coincide, proving the first assertion of \cref{thm:main}.

\subsection{Two genus checks from the plane equation}
Write $f=\sum_{j=0}^7 C_j(X)t^j$. Exact factorization gives
$C_7=aH(X)^{13}$ with $H$ an irreducible separable quartic, and
$\gcd(H,C_6)=1$. In the chart $u=1/t$, the equation
$u^7f(X,1/u)=0$ is consequently smooth at the four roots of $H$ on $u=0$,
with $u$ vanishing to order $13$. At $(X,t)=(\infty,\infty)$, put
$v=1/X$. The lowest terms of $v^{65}u^7f(1/v,1/u)$ for weights
$\mathrm{wt}(v)=7$, $\mathrm{wt}(u)=13$ are $u^7+av^{13}$.
Since $7$ and $13$ are coprime, this is one branch with
$\delta=(7-1)(13-1)/2=36$. These coefficient checks are included in
\texttt{scripts/check\_review\_evidence.py}.

The $342$ points in \cref{lem:node} are ordinary nodes: the two local roots
have distinct derivatives in $t$. All remaining points are smooth by the
preceding local checks. The closure in $\PP^1\times\PP^1$ has arithmetic
genus $(65-1)(7-1)=384$, so its geometric genus is
\[
 384-342-36=6.
\]
There is also a discriminant-degree check. The thirteen unbounded roots
grow like $t^{7/13}$; the other fifty-two approach the four distinct roots
of $H$, with differences within each group of thirteen of order
$t^{-1/13}$. Thus the squared pairwise differences give
\[
 \deg_t\disc_X(f)
 =2\binom{13}{2}\frac7{13}
  +2(13\cdot52)\frac7{13}
  -2\left(4\binom{13}{2}\right)\frac1{13}
 =84+728-48=764.
\]
Both checks agree with the certified discriminant and branch cycles.
